% YAAC Another Awesome CV LaTeX Template
%
% This template has been downloaded from:
% https://github.com/darwiin/yaac-another-awesome-cv
%
% Author:
% Christophe Roger
%
% Template license:
% CC BY-SA 4.0 (https://creativecommons.org/licenses/by-sa/4.0/)
%Section: Work Experience at the top
\sectionTitle{Educación}{\faMortarBoard}
%\renewcommand{\labelitemi}{$\bullet$}
\begin{experiences}
  \experience
    {Enero 2016} {Pontificia Universidad Javeriana}{PUJ}{Bogotá, Colombia}
    {Marzo 2022}    {
                      \begin{itemize}
                        \item \textbf{Grado en Ingeniería Electrónica}: Graduado con un promedio acumulado de 4.2 (en una escala de 0.0 a 5.0). Programa acreditado internacionalmente por ABET con formación integral en líneas clave como Sistemas de Control, Telecomunicaciones, Sistemas Digitales y Circuitos Electrónicos.
                        \item \textbf{Tesis de Grado}: Profundización y enfoque académico-práctico en el área de Procesamiento Digital y Analógico de Señales, proponiendo un método innovador basado en aprendizaje profundo para la super-resolución de imágenes plenópticas de plantas, dentro del marco del programa de investigación Optimización Multiescala In-silico de Cultivos Agrícolas Sostenibles (ÓMICAS).
                        \item Ver \href{https://wallet.xertify.co/certificates/E8381125A001}{\color{CadetBlue}{Certificado (Insignia Digital)}}
                      \end{itemize}
                    }
                    {Ingeniero Electrónico, CDIO, Procesamiento de Señales, Procesamiento de Imágenes, Visión por Computadora}

                    %{Electronic Engineer, CDIO, Signal Processing, Image Processing, Computer Vision}
\emptySeparator
  \experience
    {Septiembre 2022} {Universitat Politècnica de València}{UPV}{Valencia, España}
    {Septiembre 2023}    {
                      \begin{itemize}
                        \item \textbf{Máster en Inteligencia Artificial, Reconocimiento de Formas e Imagen Digital}: Graduado con un promedio acumulado de 8.9 (en una escala de 0.0 a 10.0). Formación avanzada y especialización en algoritmos de aprendizaje automático aplicados a Visión por Computadora (CV), Reconocimiento Automático del Habla (ASR) y Procesamiento del Lenguaje Natural (NLP).
                        \item \textbf{Tesis de Máster}: Desarrollo e implementación de un método novedoso de visión artificial para detección de objetos con estimación de profundidad integrada para sistemas en tiempo real. \href{https://riunet.upv.es/handle/10251/198268}{\color{CadetBlue}{[Explorar Tesis]}}
                        \item Ver \href{https://sede.upv.es/eVerificador}{\color{CadetBlue}{Certificado}} con el código ALUX0W631LT
                      \end{itemize}
                    }
                    {Ingeniero de Inteligencia Artificial, Aprendizaje Profundo, Visión por Computadora, Procesamiento del Lenguaje Natural}
                    %{Artificial Intelligence Engineer, Computer Vision Engineer, Machine Learning}
\emptySeparator
  \experience
    {Septiembre 2024} {Universitat Politècnica de València}{UPV}{Valencia, España}
    {Actualidad}    {
                      \begin{itemize}
                        \item \textbf{Doctorado en Informática}: Investigación doctoral centrada en algoritmos avanzados de aprendizaje automático para visión por computadora aplicados a la agricultura sostenible, en colaboración con el Valencian Research Institute for Artificial Intelligence (VRAIN).
                        \item \textbf{Mención Industrial}: Convenio industrial en colaboración con la empresa \textbf{Arisnova} por la participación activa en el proyecto de innovación \textbf{VisionTex}: Desarrollo de un Sistema de Inspección de Calidad Basado en Tecnologías de Visión Artificial para la Industria.
                      \end{itemize}
                    }
                    {Inteligencia Artificial, Visión por Computadora, Investigación, Industria}
                    %{Artificial Intelligence Engineer, Computer Vision Engineer, Machine Learning}
\end{experiences}
